% Generated from locale/tr/content/first-order-logic/introduction/first-order-logic.tex; do not hand-edit.


\olfileid[tr]{fol}{int}{fol}

\olsection{Birinci Dereceden Mantık}

Birinci dereceden mantığı büyük olasılıkla biçimsel mantıkla ilk
tanışmanızdan biliyorsunuz.\footnote{Aslında bunu aşağı yukarı
varsayıyoruz!{} Bilmiyorsanız \emph{forall x} \citep{Magnus2021} gibi
daha temel bir ders kitabını gözden geçirebilirsiniz.} Onu
``niceleme mantığı'' ya da ``yüklem mantığı'' adıyla tanıyor
olabilirsiniz. Birinci dereceden mantık her şeyden önce biçimsel bir
dildir. Yani belirli bir söz dağarcığı vardır ve ifadeleri bu söz
dağarcığından oluşturulan dizgelerdir. Fakat her dizgeye izin verilmez.
İzin verilen farklı ifade türleri vardır: terimler, !!{formula}s ve
!!{sentence}s. Biz öncelikle birinci dereceden mantığın
!!{sentence}s ifadeleriyle ilgileniyoruz. Bunlar doğal dildeki
tümcelerin biçimsel bir karşılığını verir ve onlar hakkında bir
mantıkçının genellikle ilgilendiği soruları sorabiliriz. Örneğin:
\begin{itemize}
    \item $!B$,~$!A$'dan mantıksal olarak çıkar mı?
    \item $!A$ mantıksal olarak doğru, mantıksal olarak yanlış, yoksa
olumsal mıdır?
    \item $!A$ ile $!B$ denk midir?
\end{itemize}

Bu sorular öncelikle birinci dereceden mantığın !!{sentence}s
ifadelerinin ``anlamı'' hakkındadır. Örneğin bir filozof,~$!B$'nin
$!A$'dan mantıksal olarak çıkıp çıkmadığı sorusunu şöyle çözümlerdi:
$!A$'nın doğru,~$!B$'nin yanlış olduğu bir durum var mıdır (öyleyse
$!B$,~$!A$'dan çıkmaz), yoksa $!A$'yı doğru kılan her durum~$!B$'yi
de doğru mu kılar (öyleyse $!B$,~$!A$'dan çıkar)? Fakat henüz
``durum''un ne olduğu söylenmedi; bunu belirlemek \emph{semantiğin}
görevidir. Birinci dereceden mantığın semantiği, filozofun sezgisel
``durum'' fikrine ve---bu önemlidir---bir~$!A$ !!{sentence} ifadesinin
bir durumda \emph{doğru olmasının} ne demek olduğuna matematiksel
bakımdan kesin bir model verir. Geliştireceğimiz bu matematiksel kesin
modele !!a{structure} diyoruz. ``İçinde doğrudur'' ifadesini kesin
kılan bağıntıya \emph{sağlama} bağıntısı denir. Böylece
!!{sentence}s~$!A$ ve !!{structure}s~$\Struct M$ için
``$!A$,~$\Struct M$ içinde sağlanır'' ifadesini (simgesel olarak
$\Sat M{!A}$) tanımlayacağız. Bunu yaptıktan sonra ``-dan çıkar'' ya da
``mantıksal olarak doğrudur'' gibi öteki semantik terimlere de kesin
tanımlar verebiliriz. Bu tanımlar örneğin
\[
\lforall[x][!A(x)\lif !B(x)],
\lexists[x][!A(x)] \Entails \lexists[x][!B(x)]
\]
olup olmadığını yine matematiksel kesinlikle belirlemeyi mümkün kılar.
Yanıt elbette ``evet''tir. Doğal dildeki tümceleri birinci dereceden
mantıkta simgeleştirme eğitimi aldıysanız bunu örneğin ``Bütün
karıncalar böcektir; karıncalar vardır; dolayısıyla böcekler vardır''
çıkarımının simgeleştirmesi olarak tanırsınız. Bu açıkça geçerli bir
çıkarımdır; dolayısıyla biçimsel dilimiz için kurduğumuz ``-dan çıkar''
matematiksel modeli de aynı yanıtı vermelidir.

Biçimsel mantıkla ilk tanışmanızdan muhtemelen hatırladığınız başka bir
konu da \emph{!!{derivation}s} bulunmasıdır. İlk biçimsel mantık
dersini aldıysanız eğitmeniniz bu tür !!{derivation}s bulma
alıştırmaları, belki yukarıdaki mantıksal gerektirmenin geçerli
olduğunu gösteren !!a{derivation} bile yaptırmıştır. !!{derivation}s
vermenin pek çok farklı yolu vardır: ``doğal tümdengelim'' ya da
``doğruluk ağaçları'' denen bir yöntem kullanmış olabilirsiniz, fakat
başka birçok yöntem de vardır. !!{derivation} sistemlerinin amacı,
mantıkçıların yukarıdaki sorularını yanıtlayabilecekleri araçlar
sağlamaktır. Örneğin öncülleri
$\lforall[x][!A(x)\lif !B(x)]$ ve $\lexists[x][!A(x)]$, sonucu (son
satırı) $\lexists[x][!B(x)]$ olan doğal tümdengelim
!!{derivation} ifadesi, $\lexists[x][!B(x)]$ ifadesinin ilk iki
ifadeden mantıksal olarak çıktığını \emph{doğrular}.

Peki neden? Görünüşte !!{derivation} sistemlerinin semantikle hiçbir
ilgisi yoktur: biçimsel !!{derivation} vermek yalnızca simgeleri
belirli kuralların yönettiği biçimlerde düzenlemekten ibarettir; bu
sistemler ``durumlar''dan ya da ``içinde doğru'' olmaktan hiç söz
etmez. !!{derivation} sistemleriyle semantik arasındaki bağlantının
metamantıksal bir incelemeyle kurulması gerekir. Örneğin şu matematiksel
ispat gereklidir: öncülleri $\lforall[x][!A(x)\lif !B(x)]$ ile
$\lexists[x][!A(x)]$ olan biçimsel bir
$\lexists[x][!B(x)]$ !!{derivation} ifadesi ancak ve ancak bu iki öncül
birlikte $\lexists[x][!B(x)]$ ifadesini mantıksal olarak gerektiriyorsa
mümkündür. Ancak bunu yapmadan önce sözdizim ve semantik tanımlarını
doğru kurmak için büyük ve titiz bir çalışma gerekir.

